\documentclass{article}
\usepackage{amsmath,amssymb,amsthm,hyperref}
\newtheorem{theorem}{Theorem}
\title{LSPV-3 CONTEXT-VALUE-AFFINE171}
\date{October 29, 2026}
\begin{document}\maketitle
\noindent\textbf{Status:} restricted elementary results PROVED/EXACT, independent referee pending. Greedy positive marginal-value policy FALSIFIED for explicit action menu; universal SAT and P versus NP OPEN.
\section{PROPOSE, SEARCH, BUILD}
Study the detached summary prepared on $A$ before context $B$ is chosen for $F=A(X,Y)\land B(X,Z)$. Measure construction $P$, stored bits $L$, downstream cost $Q$, peak simultaneous workspace and certificate acquisition. Storing $A$ is permitted. The joint solver seeing $A,B$ is not the same model. Authenticated sources reviewed include LSPV-3 RANK-SWITCH168, WEIGHTED-SUMMARY170, SOURCE-RUP164, LSPV-1 ARC227; actual DPMC join tree and CaDiCaL elimination source, and Darwiche--Marquis knowledge compilation and Golovin--Krause adaptive submodularity.
\begin{theorem}
For affine one-row source $A_{\alpha,\beta}: \langle\alpha,X\rangle=\beta$ with $b$ bits of $\alpha$, any deterministic fixed-length summary supporting every pin context $B_x$ has $L\ge b+1$. Copying $(\alpha,\beta)$ attains this with $P,Q=O(b)$.
\end{theorem}
\begin{proof}
The $2^{b+1}$ pairs $(\alpha,\beta)$ define distinct functions $x\mapsto[\langle\alpha,x\rangle=\beta]$: zero normals give distinct constants and nonzero normals give distinct hyperplanes. A summary collision would force an incorrect answer on some pin. Thus $2^L\ge2^{b+1}$. Original-source retention attains equality, disproving an exponential floor relative to the $b+1$-bit source.
\end{proof}
\begin{theorem}
For $b=2h$, sources $z=Mu$ with $M\in\mathbb F_2^{h\times h}$ require $L\ge h^2$ under all pins, while retaining $M$ uses $h^2+O(h\log h)$ bits.
\end{theorem}
\begin{proof}
Distinct matrices yield distinct graph predicates; pin $(u,Mu)$ for a vector where $Mu\ne M'u$ to distinguish. There are $2^{h^2}$ such functions.
\end{proof}
\section{Contextual value and false greedy guarantee}
Write $x\mid_v$ and $(x=1)\mid_v$ for contextual value labels, not probability conditioning. Use $v_\Phi(I\mid s)=\Phi(s)-\Phi(s\text{ after certified application of }I)$, with the acquisition and verification cost separately charged. For pin $X=e_i$, target $T=[\alpha_i=\beta]$ and independent uniform $\alpha_i,\beta$, $\Phi=H(T\mid K)$ gives initial individual bit gains $0,0$ but joint gain $1$. Duplicate gain is zero. Thus target information gain violates diminishing returns. For pin $X=0$, only $\beta$ matters, proving context dependence.
\begin{theorem}
If permitted actions are read $\alpha_i$ at cost 1, read $\beta$ at cost 1, or an exact direct answer macro costing $K\ge3$ charged steps, then a policy selecting maximal \emph{strictly positive individual} $v_\Phi/C$ chooses the macro and costs $K$, while two reads cost 2. The factor $K/2$ is unbounded over the action-menu parameter. This refutes the restricted greedy rule, not arbitrary SAT.
\end{theorem}
\begin{proof}Each single share has gain 0; the direct answer has gain 1, ratio $1/K>0$. Two shares determine $T$ exactly.\end{proof}
\section{Affine certificate controller}
\begin{theorem}
For affine $A:M_X X+M_Y Y=c$, compute $U$ with rows spanning $\ker(M_Y^\top)$. Its exact boundary projection is $U M_X X=Uc$, with at most $(m-\operatorname{rank}M_Y)(b+1)$ essential bits. For affine $B(X,Z)$, Gaussian elimination gives a single uniform SAT/UNSAT controller: SAT yields reconstructed $(X,Y,Z)$ checked against source; UNSAT yields an original-source indexed XOR combination $0=1$. Bit-time $O(H^4)$, simultaneous peak work bits $O(H^2)$ for combined physical encoding length $H$, charging pivot discovery, source scans, failed rows and output.
\end{theorem}
\begin{proof}
$M_Y Y=c+M_X X$ is solvable exactly if every left-kernel vector annihilates its right-hand side. Attach original-row XOR masks to each projected row. Reconstruct $Y$ for a model, or expand a Gaussian contradiction to original $A,B$ source IDs for a replayable UNSAT certificate.
\end{proof}
\section{Charged potential}
\begin{theorem}
For complete states $s_i$ with an independently defined nonnegative integer $\Phi$, initial $\Phi\le P(H)$, total upward drift $\sum_i(\Phi_{i+1}-\Phi_i)_+\le G(H)$, and at most $C(H)$ fully charged operations between successive strict drops (also initially/finally), total time is at most $(P+G+1)C$. Polynomial $P,G,C$ and same-execution peak space imply polynomial simultaneous resources.
\end{theorem}
\begin{proof}
Each strict drop consumes at least one unit. At most $P+G$ strict drops can occur, hence at most $P+G+1$ charged epochs.
\end{proof}
Affine row-processing has a polynomial-value remaining-row obligation count and polynomial work per completed row. This rank is not entropy or optimal remaining runtime. A realized entropy change can be negative; for prior $P(T=1)=.1$ and rare evidence making posterior $P(T=1)=.5$, it equals $h_2(.1)-1=-.531004$ bits while expected gain is positive.
\section{Reproducibility and next lemma}
Compact source independently checks 300 affine pairs (159 SAT, 141 UNSAT), 1402 exact projection checks and 141 original-source XOR contradictions; the full local archive tests 576 cases and 4320 projection checks (raw SHA-256 40dc5a6a07ed3c6d85d106919419920e250ffe91a5954216d2aa4d191b8b2858), but complete archive is not claimed GitHub-synced. Next: bound complete-state upward drift and charged nonprogress for connected strict exact-(2,2)-3CNF moving interfaces. P=NP and P$\ne$NP remain OPEN.
\end{document}
